\documentclass[10pt]{article}
\usepackage[letterpaper,margin=.38in]{geometry}
\usepackage{amsmath,amssymb,multicol}
\pdfinfoomitdate=1\pdftrailerid{}
\pagestyle{empty}\setlength{\parindent}{0pt}\setlength{\parskip}{2pt}
\setlength{\columnsep}{13pt}\setlength{\jot}{0pt}
\setlength{\abovedisplayskip}{3pt}\setlength{\belowdisplayskip}{3pt}
\let\bigsum\sum\renewcommand{\sum}{\bigsum\nolimits}
\newcommand{\fl}[1]{\left\lfloor#1\right\rfloor}
\newcommand{\rg}[2]{\{#1,\ldots,#2\}}
\begin{document}
\begin{center}
{\large\bfseries A complete fourteen-variable consistency equation}\\[-1pt]
{\small Strictly positive integer unknowns; exact degree 150; a 65-operation evaluator.}
\end{center}\vspace{-8pt}
\begin{multicols}{2}
\fontsize{9}{9.5}\selectfont
The only unknowns are
\[
J,F,\alpha,z,f,i,T,s,w,n,\eta,\zeta,y,\gamma\in\mathbb Z_{>0}.
\]
All other symbols in the equation are polynomial abbreviations:
\[
\begin{gathered}
b=\beta J,\quad q=b+1,\quad C=q-F-\alpha,\\
X=wq,\quad Y=sq^3,\quad E=XY,\quad k=\eta+\zeta,\\
\ell=EkY,\quad c=kY+\eta,\quad a=E+Y,\\
H=4a+3,\quad\Delta=a^2+H,\quad D=X+ac+\gamma H,\\
R=[b(q-F)+\alpha](q^2-1)+(\mu+q\omega)J,\\
\tau=En+k-R,\quad V=f^2(cT-R)-c,\\
P=\tau^2-\ell(\ell+k),\quad Q=D^2-\Delta c^2,\\
S=i\Delta c^2,\quad U=S^2(V^2-y^2)+y^2,\\
W=(K+w)C+(q-F)-zb,\quad Z=\Delta f^2-S^2.
\end{gathered}
\]
The single equation is
\[
\boxed{PQUWZ-\Delta=0.}
\]

\textbf{The four fixed integers.}
Every index, set, sum, product and floor below defines a constant, not
another unknown. Put
\[
\begin{gathered}
\rho=2^{47},\quad B=\rho^{2011},\quad\beta=B-1,\\
G_m(x)=\sum_{j=0}^{m-1}x^j,\quad
\Lambda=\rho^{1005}.
\end{gathered}
\]
\begin{align*}
\mu={}&\beta-G_{951}(\rho)\\
&-\rho^{960}G_{26}(\rho)-\rho^{1004},\\
\omega={}&\beta+8G_{917}(\rho)-4+(2^{45}-4)\Lambda\\
&+\rho^{960}G_9(\rho)+\rho^{987}G_8(\rho)
  +\rho^{1004},\\
K={}&\rho^{18}+\rho^{1004}+\rho^{2020}\\
&+2\Lambda\Phi(1/\rho),
\end{align*}
where
\[
\begin{gathered}
\chi_i=1+2^{10}(1+w_i+2^8n_i\\
\qquad+2^{17}p_i+2^{25}s_i)\quad(0\le i\le941),\\
\pi(0)=0,\quad\pi(i)=1+\sum_{j=1}^{941}
\fl{\frac{2^{44}-1+\chi_j-\chi_i}{2^{44}}}\quad(1\le i\le941).
\end{gathered}
\]
\[
\begin{split}
\Phi(x)={}&\sum_{i=0}^{941}\chi_i x^{\pi(i)}\\
&+2^{10}x^{969}(x^9+2^{17})G_8(2x)\\
&+2^{18}x^{942}(x^{18}+2^{17})G_9(2x).
\end{split}
\]
The rational intermediate in K cancels exactly; all four results are positive integers.

\textbf{The 942 indexed fixed tiles.}
Use the actions $u_j,v_j,d_j$ decoded below. For $d=0,1$, let
\[
\begin{gathered}
h(q,b)=2q+b+2,\quad
\mathcal S_d=\{u_j:d_j=d,\ 0\le j<430\},\\
r_{d,q}=205-199d+|\mathcal S_d\cap\rg0{q-1}|.
\end{gathered}
\]
An interval ending at $-1$ is empty. The table defines $t_i$ for
$0\le i\le941$; all index ranges are fixed:
\[
\setlength{\arraycolsep}{1.4pt}\begin{array}{c|l}
\text{index }i&\text{tile }t_i=(w_i,n_i,p_i,s_i)\\\hline
0&(2,435,3,1)\\
j+1\ (j=2q+b<430)&((1-d_j)r_{0,u_j},h(q,b),d_jr_{1,u_j},v_j)\\
433-2a+b\ (a,b=0,1)&(0,432a+b,0,432a+b)\\
\substack{423+2r_{d,q}+b\\q\in\mathcal S_d,\ d,b=0,1}&(dr_{d,q},b,(1-d)r_{d,q},h(q,b))\\
923+2a+b\ (a,b=0,1)&(0,432a+b,0,438)\\
927&(2,434,2,0)\\
928&(3,436,4,4)\\
929&(4,436,4,0)\\
496+t\ (t=434,435,436)&(1,438,1,t)\\
933+3a+b\ (a,b=0,1,2)&(a^2,437,b,437)
\end{array}
\]
\textbf{A closed decoder for the 430 actions.}
Put $D_b(A,k)=\fl{A/b^k}\bmod b$, where
$a\bmod b=a-b\fl{a/b}$. For $0\le q<215$, $b=0,1$, $j=2q+b$:
\[
\begin{gathered}
\mathcal N=\mathcal A\bmod2^{636},\quad
\mathcal T=\fl{\mathcal A/2^{636}},\\
\nu_j=D_2(\mathcal N,j),\quad \kappa_j=(\sigma_j+2)^{1-\nu_j}\quad(0\le j<430),\\
\sigma_j=\sum_{k=0}^{j-1}\nu_k,\quad
\Pi_j=\prod_{k=0}^{j-1}\kappa_k\quad(0\le j\le430),\\
m_q=D_2(\mathcal N,430+q),\quad \tau_q=\sum_{k=0}^{q-1}m_k,\\
F_q=14+m_q(D_{16}(\fl{\mathcal T/\Pi_{430}},\tau_q)-14),\\
u_j=(\sigma_{j+1}-(\fl{\mathcal T/\Pi_j}\bmod\kappa_j))\bmod216,\\
v_j=D_2(F_q,2b),\quad d_j=D_2(F_q,2b+1).
\end{gathered}
\]
Empty sums/products are $0/1$. Here $(v_0,d_0,u_0)=(1,1,1)$; the constructor tile
uses this actual first target, including its state renaming.

\textbf{The integer $\mathcal A$ in base 10.}
Its complete decimal expansion is the concatenation of these lines:\par
\begingroup\fontsize{7.7}{8.2}\selectfont\ttfamily
12889260410359561746736540743515845134244260640012702033212621\par
95398074180784058571518488282119861023093355389406939414042892\par
97455773992193784795705001618777025966369137954859764532832722\par
75411843410809038713502633204617333246733616940964578272360901\par
56186216508504953649192752154519268024509952082950238865468307\par
22999632482769253980274091217460707829023970817244790340704449\par
09085581271192071656170732004937337082171141082486384269072884\par
72778087419663417942892520175630649792329265834490253035577533\par
96449689808375788457488971017043689875433162972390177936123032\par
89025657757489526342497438772045714099661532481123050235898973\par
27443609903800541199351405985754278307227777584454508422717998\par
334113564409335346074892680328437\par
\endgroup
\normalfont\fontsize{7.7}{8.2}\selectfont
Lean proves that the corresponding polynomial has no positive integer zero iff
$\mathrm{Con}(\mathrm{ZFC})$, and formalizes an HA proof of the equivalence.
The printed transcription was independently checked; this page is a
specification, not the proof.
\end{multicols}
\end{document}
